\section{Conclusion and Future Work}

This work presented a systematic ablation study of DenseNet121-based multi-label chest X-ray disease classification on NIH ChestX-ray14.
We evaluated the effects of CLAHE preprocessing, CBAM attention at three insertion points, three loss functions (BCE, ASL, Focal Loss), and two optimizers (Adam and AdamW) across ten experimental configurations.

The key findings are:
\begin{itemize}
  \item CLAHE preprocessing did not consistently improve Mean Test AUC and was excluded from the final configurations.
  \item CBAM block34 (attention after DenseBlocks 3 and 4) provided the best attention placement, outperforming block4-only and block1234 variants.
  \item Removing CLAHE and using AdamW with CBAM block34 and BCE achieved the highest Mean Test AUC of \textbf{0.8371} (Experiment~08).
  \item Replacing BCE with Asymmetric Loss (Experiment~09) preserved AUC (0.8370) while substantially improving Recall (+0.099) and F1 (+0.077).
  \item Focal Loss with $\gamma{=}2.0$ did not improve over either BCE or ASL.
\end{itemize}

\textbf{Experiment~08} (DenseNet121 + CBAM block34 + BCE + AdamW) is recommended as the best AUC-oriented final model with Mean Test AUC 0.8371.
\textbf{Experiment~09} (DenseNet121 + CBAM block34 + ASL + AdamW) is recommended as the best Recall/F1-oriented alternative with Recall 0.2385 and F1 0.2736 at nearly identical AUC.

\subsection*{Future Work}

Several directions could extend this study:
\begin{itemize}
  \item \textbf{Class-specific threshold tuning}: optimizing per-class decision thresholds on the validation set could improve Precision--Recall balance for each disease.
  \item \textbf{External validation}: evaluating the models on an independent dataset (e.g., CheXpert, MIMIC-CXR) would assess generalization beyond NIH ChestX-ray14.
  \item \textbf{Model calibration}: Platt scaling or temperature calibration could align predicted probabilities with empirical event rates.
  \item \textbf{Ensemble approaches}: combining Exp~08 and Exp~09 predictions may yield improved performance on both AUC and Recall objectives.
  \item \textbf{Alternative backbones}: testing EfficientNet, ResNet-based models, or vision transformer architectures could reveal whether the observed trends generalize beyond DenseNet121.
  \item \textbf{Label noise handling}: methods designed for noisy multi-label annotations may improve performance given the known NLP-extracted label noise in ChestX-ray14.
  \item \textbf{Localization evaluation}: if bounding box annotations are available (e.g., from the NIH localization subset), formal localization accuracy could be evaluated alongside Grad-CAM qualitative observations.
\end{itemize}
